\section{Essential incompleteness of arithmetic}\label{sec:arithmetic}

For $e\in\mathcal I(X,Y)$, define the partial function
$\varphi_e:|X|\rightharpoonup|Y|$ by
\[
 \varphi_e(x)=y\quad\Longleftrightarrow\quad
 \N\models G_e\bigl(\gamma_X(x),\gamma_Y(y)\bigr).
\]
Let $\mathrm{Rec}$ have the terms in $\mathsf O$ as objects and
these partial functions as morphisms, with extensional equality.
The numerical indexing represents every partial recursive function,
so this is a presentation of the recursion category of partial recursive
functions on these objects.
Its restriction idempotents are partial identities on c.e.\ subsets;
the complemented ones are partial identities on recursive subsets.

Let $u$ be the evaluator index of Lemma~\ref{lem:arithmetized-evaluation}.
Its standard denotation
$\tau_{\mathrm{rec}}=\varphi_u:\omega\times\omega\rightharpoonup\omega$
is a Turing morphism in $\mathrm{Rec}$ by the specialization identity
of that lemma. Let $e_d$ be the index obtained by composing $u$
with the chosen diagonal index. Thus
$\varphi_{e_d}=\tau_{\mathrm{rec}}\Delta_\omega$.
Choose the coproduct isomorphism $\zeta:\omega\to\omega+\omega$
whose underlying function sends $2m$ to the first tagged copy of $m$
and $2m+1$ to the second. Define
\[
 \alpha(x)\equiv\exists m\,G_{e_d}(x,2m),\qquad
 \beta(x)\equiv\exists m\,G_{e_d}(x,2m+1).
\]
Their standard truth sets
\[
 A=\{n:\N\models\alpha(\bar n)\},\qquad
 B=\{n:\N\models\beta(\bar n)\}
\]
are the even-output and odd-output sets for self-application.
Their recursive inseparability is another instance of
Theorem~\ref{thm:local-lifting}. To see this, let
$J_{\mathrm{rec}}:\mathrm{Rec}\to\Pfn$ be given by
$J_{\mathrm{rec}}(X)=|X|$ and $J_{\mathrm{rec}}(f)=f$.
It is a covariant range functor. The category $\Pfn$ satisfies
the target hypotheses: partial identities on singletons detect partial
functions, restriction idempotents split through their subsets of
definition, partial injections have partial inverses, and inverse images
preserve unions. Complementary restriction idempotents on an object $W$
in $\mathrm{Rec}$ are partial identities on complementary recursive
subsets $C,C'$ of $J_{\mathrm{rec}}(W)$. The morphism $s:W\to W+W$
given numerically by $s(x)=i_1x$ for $x\in C$ and $s(x)=i_0x$
for $x\in C'$ is total recursive. Its image under $J_{\mathrm{rec}}$ is the
required section. This verifies the lifting hypothesis of
Theorem~\ref{thm:local-lifting} for $J_{\mathrm{rec}}$.
Applied with $\tau_{\mathrm{rec}}$ and $\zeta$, the theorem gives
the partial identities on $A$ and $B$ as its two inseparable
restriction idempotents. Since complemented restriction idempotents
in $\mathrm{Rec}$ are partial identities on recursive subsets,
$A,B$ are recursively inseparable.
For $n\in A$, its finite computation trace gives a PA-proof of
$\alpha(\bar n)$; for $n\in B$, it gives a PA-proof of
$\beta(\bar n)$. PA-provable functionality of $G_{e_d}$ gives
$\PA\vdash\forall x\,\neg(\alpha(x)\land\beta(x))$.

The logical objects used below are constructed in the total part of
the restriction splitting. Let $\mathcal C_{\mathrm{rec}}$ have c.e.\
subsets of finite powers of $\N$ as objects and functions computed
by partial recursive programs total on their source carriers as
morphisms, with extensional equality.

\begin{proposition}\label{prop:arithmetic-theory-realization}
There is an equivalence
\[
 \operatorname{Tot}(\operatorname{Split}(\mathrm{Rec}))
       \simeq\mathcal C_{\mathrm{rec}}.
\]
The category $\mathcal C_{\mathrm{rec}}$ is regular, with finite
coproducts that are disjoint and stable under pullback. Every
monomorphism is an admissible domain monomorphism, and
\[
 \operatorname{Par}(\mathcal C_{\mathrm{rec}},\operatorname{Mono})
       \simeq\operatorname{Split}(\mathrm{Rec}).
\]
For $V=\N$ and a recursive signature of bounded finite arity, the
syntax and list algebras of Definition~\ref{def:categorical-syntax}
exist in $\mathcal C_{\mathrm{rec}}$ and are stable under parameters.
For every axiom subobject of the sentence object, the inference
endofunctor of Theorem~\ref{thm:first-order-theory} has an initial
algebra of finite derivations.
\end{proposition}
\begin{proof}
A restriction idempotent in $\mathrm{Rec}$ is the partial identity
on a c.e.\ carrier. A total morphism between two split objects is a
partial recursive function with exactly the source carrier as domain
and values in the target carrier. Conversely, a program total on a
c.e.\ source can first semidecide membership in that source, giving
the same function with exactly that domain. The recursive numerical
codings identify these carriers with c.e.\ subsets of finite powers
of $\N$ and give the stated equivalence.

A singleton is terminal. Products and equalizers are c.e.\ carriers,
computed by pairing and numerical equality. Tagged unions give
finite coproducts. For a morphism into a tagged union, the two tag
tests decompose its source into c.e.\ subcarriers; their coproduct
maps isomorphically to the source. Thus coproducts are disjoint and
stable under pullback. The image of a morphism is c.e.\ by
dovetailing its computations; factorization through that image is
a surjection followed by an inclusion. A surjection $e:C\to D$ is the
coequalizer of its kernel pair: for $h:C\to Z$ constant on its
fibers, search on input $y\in D$ for $x$ with $e(x)=y$ and return
$h(x)$. This is the unique descended morphism. Pullbacks preserve
surjections, proving regularity.

Singleton maps detect noninjectivity, so monomorphisms are exactly
the injections. Preimage search for an injection $m:C\to D$ gives
a partial recursive inverse on its c.e.\ image. Hence every such
monomorphism is an admissible domain inclusion up to isomorphism.
The span $C\xleftarrow{m}B\xrightarrow{g}D$ represents the partial
function $gm^{-1}$. Conversely, restricting a partial recursive
function to its domain gives such a span. These assignments respect
composition by pullback and give the restriction equivalence.

Numerical equality makes $V=\N$ decidable. Encode terms and formulas
as finite labeled ordered trees, and lists as finite sequences of
variable names. Their carriers and constructors are recursive.
To prove initiality with parameters, fix any c.e.\ carrier $P$ and
an algebra for the functor $K_P$ of Appendix~\ref{app:syntax-theories}
on $x:X\to P$ and $y:Y\to P$. For $(p,t)$ or
$(p,\varphi)$, evaluate the tree from its leaves to its root using
that algebra, retaining $p$ throughout. The algebra operations are
recursive on their specified carriers and preserve their projection
to $P$. Induction on the tree therefore proves termination and
membership of the result in the fiber over $p$. The resulting
morphisms $P\times\mathsf{Tm}\to X$ and
$P\times\mathsf{Fm}\to Y$ satisfy the algebra equations;
the same induction proves uniqueness.
This proves initiality in every such slice. Evaluation of finite
lists proves parameter stability of the list algebra in the same
way. The free-variable and substitution folds of the appendix
therefore exist here. Induction identifies them with the lists of
free-variable occurrences and the usual free-for substitution on
these named-variable trees.

Represent an axiom subobject by its c.e.\ image in the sentence
carrier. A finite derivation tree records its rule instance and
ordered children at each node. Tree shape, logical instance
conditions, and equality of child conclusions with the prescribed
premises are recursive tests. The finitely many nonlogical axiom
membership tests are semidecidable. Running these tests in parallel
gives a c.e.\ carrier of valid derivations. Reading its conclusion
is recursive, as are the operations adjoining an axiom leaf or an
inference node. In any algebra for the inference endofunctor,
evaluate a valid derivation from its leaves to its root. Induction
shows that each intermediate result has the required conclusion,
so every algebra operation is applied on its domain, and evaluation
terminates. This defines an algebra morphism; induction on the same
trees proves its uniqueness. Thus the derivation algebra is initial.
\end{proof}

Work in $\operatorname{Split}(\mathrm{Rec})$.
For a c.e.\ subset $D\subseteq\N$, write $D$ for the split object
$(\omega,1_D)$, where $1_D$ is the partial identity on $D$.
A partial recursive function $h:\N\rightharpoonup\N$ with domain
contained in $D$ and image contained in a c.e.\ subset $E\subseteq\N$
represents a morphism
$D\rightharpoonup E$; it is total precisely when its domain is $D$.
Proposition~\ref{prop:arithmetic-theory-realization} identifies this
presentation with partial maps between c.e.\ numerical carriers.
In this presentation, restriction idempotents are partial identities
on c.e.\ subsets, their joins are partial identities on unions, and inverse images along
total morphisms preserve these joins.
The full restriction functor
$I:\mathrm{Rec}\to\operatorname{Split}(\mathrm{Rec})$ is given by
$I(X)=(X,1_X)$ and $I(f)=f$. It leaves the restriction idempotents
on each original object, their order, and their complements unchanged.
Thus the partial identities $a,b$ on $A,B$ remain inseparable
on $I(\omega)$.

Use the syntax of Proposition~\ref{prop:arithmetic-theory-realization}
for the arithmetic signature. Its sentence object $S$ is the recursive
set of codes of closed arithmetic sentences.
Negation and numeral substitution give total maps
\[
 \nu:S\to S,\qquad \nu(\ulcorner\varphi\urcorner)
     =\ulcorner\neg\varphi\urcorner,
 \qquad
 \theta:\N\to S,\qquad \theta(n)=\ulcorner\alpha(\bar n)\urcorner.
\]
For a c.e.\ theory $T\supseteq\PA$, use the Hilbert calculus of
Appendix~\ref{app:first-order-derivations}. Let $\mathrm{Prf}_T$ be the
set of codes of finite derivation trees from $T$ whose conclusions are sentences.
This carrier is c.e.: the logical rules and formula syntax are
recursive, and each candidate proof has only finitely many axiom
membership tests in the c.e.\ axiom set of $T$. Reading the root
conclusion gives a total morphism
$c_T:\mathrm{Prf}_T\to S$ in the splitting. Its range
$t_T=\ran(c_T)$ is the partial identity on the c.e.\ set of theorems
of $T$.
Proposition~\ref{prop:arithmetic-theory-realization} supplies the initial
derivation algebra. Theorem~\ref{thm:first-order-theory} identifies its
sentence-conclusion pullback with $c_T$ and its image with the least
sentence theory containing the axioms of $T$.

The finite computation proofs of $\alpha(\bar n)$ for $n\in A$ give
$a\le\theta^*t_{\PA}$. For $n\in B$, a finite computation proves
$\beta(\bar n)$ in $\PA$; the PA-provable disjointness of $\alpha$
and $\beta$ then proves $\neg\alpha(\bar n)$. Thus
\begin{equation}\label{eq:arithmetic-domain-inclusions}
 a\le\theta^*t_{\PA},\qquad
 b\le\theta^*\nu^*t_{\PA},\qquad
 t_{\PA}\le t_T.
\end{equation}

\begin{corollary}[Essential incompleteness of arithmetic]\label{cor:arithmetic}
The theorem and refutation sets of every consistent computably
enumerable extension of $\PA$ are recursively inseparable.
Consequently every such extension is incomplete.
\end{corollary}
\begin{proof}
Let $T\supseteq\PA$ be consistent and c.e.\ Consistency says that
no sentence and its negation are both theorems, so
$t_T\,\nu^*t_T=0_S$. Completeness would say that every sentence or
its negation is a theorem, giving $t_T\vee\nu^*t_T=1_S$.
The pair $a,b$ supplied by Theorem~\ref{thm:local-lifting}, together
with \eqref{eq:arithmetic-domain-inclusions}, satisfies the hypotheses
of Corollary~\ref{cor:domain-incompleteness}. That corollary makes
$t_T,\nu^*t_T$ inseparable and rules out the completeness equation.
A complemented restriction idempotent on the recursive sentence carrier
$S$ is a partial identity on a recursive subset of $S$, so this gives
recursive inseparability of the theorem and refutation sets.
\end{proof}

\subsection{Numeral instances and theory extension}

For a consistent c.e.\ extension $T$ of $\PA$, define
\[
 P_T=\{n:T\vdash\alpha(\bar n)\},\qquad
 N_T=\{n:T\vdash\neg\alpha(\bar n)\}.
\]
These are the domains of definition of the partial identities
$\theta^*t_T$ and $\theta^*\nu^*t_T$.
They are computably enumerable disjoint sets with $A\subseteq P_T$
and $B\subseteq N_T$, and they enlarge under theory extension.
Let $\operatorname{Pair}(\mathrm{Rec})$ have objects $(X,A,B)$,
where $X\in\mathsf O$ and $A,B\subseteq J_{\mathrm{rec}}(X)$
are disjoint computably enumerable subsets.
A morphism $(X,A,B)\to(Y,P,N)$ is a total morphism $f:X\to Y$
of $\mathrm{Rec}$ satisfying $J_{\mathrm{rec}}(f)(A)\subseteq P$ and
$J_{\mathrm{rec}}(f)(B)\subseteq N$. Identities and composition are inherited
from $\mathrm{Rec}$. The assignment
$T\mapsto(\omega,P_T,N_T)$, with the identity function for theory
inclusions, is a covariant functor from the poset of consistent
computably enumerable extensions of $\PA$ to
$\operatorname{Pair}(\mathrm{Rec})$.
The inclusions $A\subseteq P_T$ and $B\subseteq N_T$ give the morphism
$1_\omega:(\omega,A,B)\to(\omega,P_T,N_T)$ in this category.
A recursive separator pulls back along a total recursive function;
consequently $P_T,N_T$ are recursively inseparable.


For a restriction idempotent
$c=[C]_T$ on $\omega$, define the subset of the standard natural numbers
\begin{equation}\label{eq:observation}
 E_T(c)=\{n\in\N:T\vdash C(\bar n)\}.
\end{equation}
It is computably enumerable by proof enumeration and is independent of
the representative. The operation is monotone, preserves finite meets,
and preserves top and bottom. Choose a constant-function index
$\operatorname{const}(n)\in\N$ for which
$\PA\vdash\forall x\forall y\,
(G_{\operatorname{const}(n)}(x,y)\leftrightarrow y=\bar n)$, and put
$\kappa_n=[\operatorname{const}(n)]_T:\omega\to\omega$. Then
\begin{equation}\label{eq:faithful-observation}
 n\in E_T(c)
 \quad\Longleftrightarrow\quad c\kappa_n=\kappa_n
 \quad\Longleftrightarrow\quad
 F_T(c)F_T(\kappa_n)=F_T(\kappa_n).
\end{equation}
The composition formula makes
$G_{\operatorname{comp}(j_C,\operatorname{const}(n))}(x,y)$
T-provably equivalent to $C(\bar n)\land y=\bar n$, which gives
the first equivalence.
The second follows from faithfulness of $F_T$.


For $T=\operatorname{Th}(\N)$, the identity-on-objects assignment
\[
 D:S_{\operatorname{Th}(\N)}\longrightarrow\mathrm{Rec},\qquad
 D([e]_{\operatorname{Th}(\N)})=\varphi_e,
\]
is an isomorphism: two classes are equal exactly when their uniform
graph-equivalence sentence is true in $\N$, which is equivalent
to equality of their standard partial functions.
The functor $F_T$ still exists, and $S_T$ has inseparable pairs,
although $T$ is complete.
Thus inseparability in $S_T$ alone does not imply incompleteness of $T$.

Bernardi's numerical argument uses a precomplete equivalence
relation $R$ on $\N$: for every partial recursive function $f$
there is a total recursive function $g$ such that
$f(n)\mathord\downarrow$ implies $f(n)\,R\,g(n)$.
For such a relation, distinct equivalence classes are recursively
inseparable; consequently the same holds for disjoint nonempty
unions of classes. The relation $\equiv_{\PA}$ on indices defined
in \eqref{eq:uniform} is one of his examples
\cite[pp.~125--126]{Bernardi}.
The argument above obtains $A,B$ from
Theorem~\ref{thm:local-lifting} and uses finite PA-proofs to
establish $A\subseteq P_T$ and $B\subseteq N_T$.
